\section{Related Work}
\label{sec:related}

\subsection{Frequency Bias and Subsampling in Word2Vec (2013)}

The earliest large-scale embedding systems already encountered statistical
dominance of frequent tokens.
In the original Word2Vec formulation~\cite{mikolov2013efficient},
high-frequency words were observed to contribute disproportionately
to training updates.
To mitigate this, subsampling was introduced:
frequent words are randomly discarded with probability increasing
in their corpus frequency.

The stated motivation was computational efficiency.
However, from a geometric perspective,
subsampling implicitly suppresses a dominant first-order component
associated with corpus frequency.
This can be interpreted as an early heuristic attempt to control
frequency-driven anisotropy in the embedding space.

Importantly, Word2Vec did not explicitly analyze the embedding mean
or its contribution to spectral structure.
Bias mitigation was treated as an optimization trick,
rather than as a structural property of the learned representation.

\subsection{Common Direction Removal in Sentence Embeddings (SIF, 2017)}

Arora et al.~\cite{arora2017simple} proposed Smooth Inverse Frequency (SIF)
sentence embeddings, observing that averaged word vectors tend to
be dominated by a common direction.
They remove the top principal component after sentence averaging,
demonstrating improved downstream performance.

Their interpretation is that frequent words induce a dominant
latent direction that does not encode semantic content.
This work introduced an explicit spectral removal step
(post-hoc PCA subtraction) and formalized a generative model
in which a common discourse vector explains the phenomenon.

While SIF identified and removed a dominant direction,
it did not distinguish between mean-induced structure
and higher-order covariance anisotropy,
nor did it analyze the impact on training dynamics or optimization.

\subsection{Anisotropy in Word Embeddings (All-but-the-Top, 2018)}

Mu and Viswanath~\cite{mu2018all} systematically analyzed anisotropy
in pretrained word embeddings.
They observed that embeddings concentrate in a narrow cone
and that removing the mean and the top principal components
significantly improves isotropy and performance.

The All-but-the-Top (ABTT) procedure subtracts the mean
and removes several leading singular vectors.
This explicitly connects anisotropy to low-rank structure
and demonstrates that dominant spectral directions
can be detrimental to representation quality.

However, ABTT operates as a post-processing technique
on static embeddings.
It does not analyze the causal origin of the dominant component,
nor does it connect mean-dominated structure to
optimization instability or quantization pathologies.

\subsection{Positioning of the Present Work}

The prior literature establishes three important facts:

\begin{itemize}
\item High-frequency tokens induce dominant directions.
\item Embedding spaces are strongly anisotropic.
\item Removing leading spectral components improves representation quality.
\end{itemize}

Our work differs in scope and focus.

First, we identify that in modern transformer training,
the dominant first-order structure is specifically an
\emph{activation mean bias} emerging from embedding statistics
and token frequency imbalance.

Second, we show that this bias is not merely a representation artifact,
but propagates multiplicatively through gradient outer products,
inflating optimizer second moments and quantization scales.

Third, we provide a unifying causal analysis linking
frequency imbalance $\rightarrow$ embedding bias
$\rightarrow$ activation anisotropy
$\rightarrow$ gradient-scale distortion.

Thus, while earlier work removed dominant directions
to improve embedding quality,
we show that bias is a central structural driver
of multiple downstream training phenomena
in large language models.:w
